\section{Discussion}
\label{sec:discussion}

\subsection{Principal Findings}
SAFE-3D extends camera-specific construction monitoring into a common site-coordinate representation by combining a UAV-derived 3D map with multiple fixed CCTV cameras. The results support three main findings. First, the UAV map provides a geometric reference that allows observations from otherwise incomparable camera views to be interpreted in one metric frame. Second, the localization ablation shows that a representative ground-contact anchor, metric depth-scale correction, and covariance-aware fusion contribute sequentially to geometric consistency. Third, the reconstructed trajectories support physically interpretable spatial outputs, including worker--worker separation, worker--equipment relations, and distances to predefined hazard regions.

The validation results should be interpreted at the appropriate level. Dataset A demonstrates strong internal cross-view consistency and isolates the contribution of each localization component. Dataset B provides the principal independent metric validation through leave-one-out evaluation against total-station measurements. Dataset C demonstrates field-level identity continuity and trajectory reconstruction but does not establish independent absolute point accuracy. Dataset D isolates the association module under an alternative ground mapping and is therefore auxiliary rather than an end-to-end validation of the UAV-referenced pipeline.

\subsection{Implications for Construction Safety Monitoring}
Vision-based construction-safety research has predominantly advanced image-space recognition of PPE, unsafe behavior, and interactions. However, collision, intrusion, and equipment-approach risks are governed by physical relationships rather than visual separation alone. SAFE-3D reinterprets existing CCTV cameras as spatially aware sensors by registering them to a metric site frame. This connection is particularly relevant because UAV mapping, as-built documentation, progress monitoring, and digital-twin workflows already generate static spatial representations, whereas CCTV provides persistent observations of dynamic site actors. SAFE-3D links these static and dynamic information sources in one coordinate system.

The proximity analysis quantifies why this connection matters. Pixel-distance discrimination varied substantially between cameras and produced false-close cases caused by depth ambiguity. A single image-plane threshold therefore cannot provide a camera-independent safety distance. Metric 3D trajectories do not automatically constitute a complete warning system, but they provide the spatial substrate required to define equipment envelopes, restricted regions, and distance-based rules in physical units.

\subsection{Role of Ground-Contact Localization and Uncertainty}
The box center does not represent the ground position of a worker or item of equipment. The raw bottom-center is closer to the expected contact region but remains sensitive to bounding-box jitter, shadows, truncation, and partial lower-body occlusion. The proposed GCA is therefore described conservatively as a representative ground anchor rather than a directly detected footpoint. Its role is to reduce systematic image-anchor inconsistency before depth-based lifting.

Scale correction is equally important because relative monocular depth cannot be interpreted directly in meters. Aligning depth-derived geometry to the metric site frame produced the largest reduction after anchor correction in the ablation study. The final fusion stage further reduced reprojection error by weighting camera observations according to their uncertainty.

The uncertainty model is most valuable when observations differ in quality. A nearby frontal observation, a distant oblique observation, and a partially occluded observation should not contribute equally to fusion or association. Observation covariance allows SAFE-3D to weight these cases differently and to interpret spatial discrepancy relative to predicted uncertainty. However, the present covariance should be viewed as a geometry-informed approximation. A stronger future evaluation should examine calibration using empirical coverage, normalized innovation statistics, or error stratification by distance and viewing direction.

\subsection{Cross-Camera Association and Threshold Transfer}
The association results show that a fixed Euclidean radius is sensitive to site geometry. A radius that is sufficiently large to retain distant uncertain observations can merge nearby objects, whereas a smaller radius can reject correct matches under long viewing distances. The GT-tuned Euclidean result demonstrates that distance-based association can perform well when the evaluation identities are available for tuning, but this condition is not representative of a newly installed site.

SAFE-3D does not eliminate all parameters. It replaces scene-specific Euclidean radius tuning with a statistical compatibility rule based on propagated track covariance and observation covariance. The confidence level, process noise, confirmation length, and retention settings remain fixed system parameters. This distinction is important: the contribution is cross-site reuse of a common uncertainty-aware association rule, not parameter-free tracking.

The re-linking stage addresses a complementary failure mode. Conservative online gating favors track fragmentation over incorrect merging during long observation gaps. Motion-gated re-linking recovers compatible tracklets afterward while retaining the same covariance-normalized criterion. It should therefore be interpreted as a continuity-recovery mechanism rather than an independent tracking contribution. A complete final table should report IDF1, HOTA, Pair-F1, IDSW, and reconstructed identity count from the same latest run to quantify this trade-off without mixing experimental configurations.

\subsection{Practical Deployment Considerations}
The system naturally separates into offline and online stages. UAV reconstruction, metric alignment, and CCTV registration can be performed offline, whereas detection, lifting, association, fusion, and trajectory update operate on incoming camera streams. Construction-site geometry is not static, however. Excavation, backfilling, structural growth, and temporary facilities can invalidate parts of the map. Practical deployment therefore requires a map-update policy based on periodic UAV acquisition or change detection.

CCTV pose accuracy directly limits localization accuracy. Cameras displaced by vibration, impact, or maintenance require re-registration. Long-term deployment would benefit from automatic pose-drift detection using repeated reference matching. Operational visualization should expose global tracks, camera support, uncertainty, and zone distance on a site or BEV map so that managers can interpret why a spatial relation was generated.

Downstream warning logic should combine distance with equipment type, operating radius, speed, work-zone rules, and uncertainty. Near a decision threshold, a probabilistic or interval-based policy is more defensible than a deterministic threshold applied to a single estimated distance. The current study establishes the geometry and trajectory representation needed for such rules but does not claim to validate a complete hazard-warning policy.

\subsection{Limitations and Future Work}
Several limitations remain. First, localization accuracy depends on CCTV calibration, map quality, depth estimation, and GCA reliability. Distant workers, prolonged occlusion, and incomplete lower-body observations increase uncertainty and may require semantic ground constraints, stronger temporal smoothing, or explicit contact-point estimation. Second, the operating-site dataset lacks independent point-wise total-station validation because continuous target installation and surveying were not feasible during active operations. Future field studies should use non-interfering survey protocols, robotic total stations, or synchronized reference sensors.

Third, cross-camera association remains sensitive to missed detections, severe occlusion, synchronization error, and dense interactions. Geometry and motion can be complemented by short-term appearance evidence, although visually similar PPE limits its role. Fourth, the four datasets cover complementary conditions but remain limited in site type, camera height, equipment diversity, density, illumination, and construction phase. Broader multi-site validation is required. Fifth, the map-maintenance interval and the effect of progressive site change were not evaluated. Future work should investigate dynamic map updates, online pose refinement, calibrated uncertainty, uncertainty-aware risk rules, and integration with BIM and construction digital twins.
